
%\needspace{4\baselineskip}
\color{uniblue}{\subsection{Контроль отсутствия напряжения на шинах (КОНш)}}
\color{black}

\begin{enumerate}

\item
В устройстве контроль отсутствия напряжения на шинах представлен функцией «КОНш». Функция контролирует отсутствие напряжения на секции шин. Функция используется для АПВ, а также для включения с контролем синхронизма. 
% УСЛОВНОЕ ВЫРАЖЕНИЕ
\ifboolexpr{bool {isVV}}
{
    Функция используется для АВР, ВНР и АПВ с контролем напряжения на секции шин, а также для оперативного включения.
}
{
    Функция используется для АПВ с контролем напряжения на секции шин, а также для оперативного включения.
}

\item
Пусковые условия функции формируются без выдержки времени при снижении всех линейных напряжений ниже уставки и одновременном выполнении следующих условий если уставка SGF2 выставлена в положение «Измеренный»:
\begin{itemize}
    \item функция введена в работу;
    \item минимальное из подведенных линейных напряжений, измеренных с шин ниже уставки;
    \item отсутствует срабатывание ИО напряжения обратной последовательности;
    \item неактивный сигнал срабатывания БННш.
\end{itemize}

\item
Предусмотрена блокировка функции при срабатывании ИО обратной последовательности напряжения выше уставки и наличии сигнала срабатывания БННш.

\item
Если уставка SGF2 выставлена в положение «Внешний» работа функции осуществляется при наличии активного внешнего дискретного сигнала «Наличие Uш» без учета срабатывания измерительных органов.

\item
Ввод функции в работу на этапе параметрирования устройства осуществляется программным переключателем SGF1 «Ввод\_функции». Информация о введенной функции отображается выходным сигналом «КОНш: Ввод».

\item
Функция может быть выведена из работы оперативно путем активации сигнала «ОВ КОНш», что характеризуется активным состоянием сигнала «КОНш: Оперативный вывод» на выходе функции.

\item
Алгоритм работы функции выполнен в соответствии с рисунком \ref{konsh:img1}. В таблице \ref{konsh:tbl1} приведены параметры, необходимые для настройки функции.

\vspace{1mm}
\begin{figure}[H]
\centering
\includegraphics[width=0.9\textwidth,height=0.9\textheight,keepaspectratio]{\fbpath/124. КОНш/_latex/img/konsh.pdf}
\captionof{figure}{Функциональная схема логики КОНш}\label{konsh:img1}
\end{figure}

\end{enumerate}
